\documentclass{article}
\usepackage[T1]{fontenc}
\usepackage{amsmath}
\usepackage{listings}
\usepackage{xcolor}

\begin{document}

star map generator notes, mostly for nothing- i just wanna type
\bigskip

\section{theory}

ok so the whole project is basically one question dressed up as a python script: given a star's fixed spot in the sky, and where i am + what time it is, where is that star right now.

so there's two coordinate systems having a little disagreement here.

\smallskip
\textbf{equatorial coordinates (RA/dec)} — this is how the catalog stores stars, basically lat/long but for the sky. 

RA is like longitude except measured in hours instead of degrees because apparently old astronomers looked at 360$^\circ$, a perfectly normal number and said non merci, we're doing clocks instead. 

declination is like latitude, $-90^\circ$ to $+90^\circ$, much more reasonable behaviour. since stars are stupidly far away, these barely budge over a human lifetime — so one RA/dec pair per star, forever et catalog done.
\medskip

since 24h = 360$^\circ$,
\begin{equation}
1^{\text{h}} = 15^\circ
\end{equation}
so converting RA hours $\to$ degrees is just times 15. i will forget this by tuesday, hence it is now written down forever.

\bigskip
\textbf{horizontal coordinates (alt/az)} — this one i actually care about because it gets plotted. altitude = how high up (0$^\circ$ = horizon, 90$^\circ$ = directly overhead, negative = below the horizon = don't bother looking, it's not there). azimuth = compass direction, 0$^\circ$ north, 90$^\circ$ east blah blah.

annoying part: this depends on where you're standing AND what time it is, because earth is rude and keeps spinning without asking anyone.

\end{document}

